\clearpage
\documentclass[../../main.tex]{subfiles}

\begin{document}

\begin{figure}[htbp]
\centering

\begin{tcolorbox}[
    title=\textbf{Error Type - Format Error},
    colback=orange!1!white,
    colframe=orange!50!black,
    fonttitle=\bfseries,
    enhanced jigsaw,
    breakable
]
\small


\begin{tabularx}{\linewidth}{@{}lX@{}}
\textbf{task type} & alfworld \\
\textbf{task description} & clean some plate and put it in shelf. \\
\textbf{critical failure module} & action \\
\textbf{evidence} &Step 15 has a non-empty plan but the recorded action is empty. Step 16 observation returns 'Nothing happens'. From Step 17 onward, admissible actions at the dining table never include 'take plate 1 ...', yet attempts persist. \\
\textbf{critical step} & 15 \\
\end{tabularx}

\medskip
\textbf{step trajectory}

\begin{Verbatim}[breaklines=true, breakanywhere=true]
observations: "On the diningtable 1, you see a apple 1, a bread 1, a butterknife 1, a cup 3, a fork 1, a houseplant 1, a knife 1, a lettuce 3, a lettuce 2, a peppershaker 2, a plate 2, a plate 1, a saltshaker 1, a soapbottle 2, a soapbottle 1, a spoon 2, a spoon 1, a winebottle 3, and a winebottle 2."

admissible_actions: "'examine diningtable 1'\n 'examine dishsponge 1'\n 'go to cabinet 1'\n 'go to cabinet 10'\n 'go to cabinet 11'\n 'go to cabinet 12'\n 'go to cabinet 13'\n 'go to cabinet 2'\n 'go to cabinet 3'\n 'go to cabinet 4'\n 'go to cabinet 5'\n 'go to cabinet 6'\n 'go to cabinet 7'\n 'go to cabinet 8'\n 'go to cabinet 9'\n 'go to coffeemachine 1'\n 'go to countertop 1'\n 'go to countertop 2'\n 'go to drawer 1'\n ..."

memory: "Looking at the past observations and actions from our conversation history, a relevant memory for this step is from step 10, where the observation was \"You open the drawer 1. The drawer 1 is open. In it, you see a dishsponge 1.\" and the action was \"take dishsponge 1 from drawer 1\". This memory is helpful because it shows that the dishsponge 1, which is likely needed for cleaning, has been retrieved and is now in the agent's possession. Another relevant memory is from step 14, where the observation was ..."

reflection: "Reflecting on the last action and its outcome, the last action was \"examine diningtable 1\" and the outcome was that the agent observed the current state of the diningtable 1. The task goal is to clean some plate and put it in a shelf, and the last action did not directly contribute to completing this goal. However, it did provide information about the current state of the diningtable 1, which is necessary for planning the next step. The agent is making progress toward the task goal by gathering information, but the actual cleaning and putting away of the plate has not started yet. To make progress, the agent should use the dishsponge 1 to clean a plate and then put it away in a shelf.",
      
plan: "Planning the next step based on memory and reflection, the agent should use the dishsponge 1 to clean a plate. Since the agent is currently at the diningtable 1 and there are plates present, the agent should take a plate from the diningtable 1 and clean it using the dishsponge 1. After cleaning the plate, the agent can put it away in a shelf. This plan is helpful because it directly addresses the task goal and uses the information gathered from previous actions. The expected outcome of this action is that the agent will have a clean plate that can be put away in a shelf."

action: "take plate 1 from diningtable 1."
\end{Verbatim}

% \vspace{2mm}

\end{tcolorbox}
\vspace{-1em}
\caption{\textbf{Error Type: Format Error}}
\label{fig:error_format_error}
\end{figure}

\end{document}